\chapter{Skin Tags (Acrochordons)}

\section*{Definition}
Skin tags (acrochordons) are small, soft, benign skin growths that hang off the skin by a thin stalk (peduncle). They are very common, affecting up to 60 percent of adults, and increase in frequency with age.

\section*{What Happens in Your Body}
Skin tags are composed of loose fibrous tissue and are covered by epidermis. They develop in areas of skin friction, such as the neck, armpits, groin, and eyelids. Hormonal changes (pregnancy, growth hormone), insulin resistance, and genetic factors contribute to their formation. Human papillomavirus (HPV) has been implicated in some cases.

\section*{Causes}
\begin{itemize}
\item Friction from skin rubbing against skin or clothing
\item Obesity (more skin folds and friction)
\item Aging (increased prevalence after age 40)
\item Pregnancy (hormonal changes)
\item Insulin resistance and type 2 diabetes
\item Acromegaly
\item Genetics: family history
\item HPV infection (some evidence)
\end{itemize}

\section*{When to See a Doctor}
\begin{itemize}
\item Small, soft, flesh-colored or slightly darker growth
\item Pedunculated (attached by a stalk)
\item Typically 2-5 mm but can be larger
\item Most common on the neck, axillae, groin, and eyelids
\item Usually painless but can become irritated from friction
\item May twist (torsion) causing pain and discoloration
\item Cosmetic concern is the most common reason for removal
\end{itemize}

\section*{Related Symptoms}
\begin{itemize}
\item Skin growths and lesions
\item Moles (nevi)
\item Warts (verrucae)
\item Seborrheic keratosis
\item Neurofibromas
\end{itemize}

\section*{Self-Care / Home Management}
\begin{itemize}
\item Do NOT attempt to cut or tie off skin tags at home
\item Keep the area clean and dry
\item Avoid irritation from clothing or jewelry
\item Monitor for changes in size, shape, or color
\item Consult a healthcare provider for removal if bothersome
\end{itemize}

\section*{How Your Doctor Will Evaluate You}
\begin{itemize}
\item Clinical examination is usually sufficient
\item Dermoscopy for confirmation
\item Biopsy if atypical features suggest other pathology
\end{itemize}

\section*{Treatment Options}
\begin{itemize}
\item Snip excision with scissors (with or without local anesthesia)
\item Cryotherapy (liquid nitrogen freezing)
\item Electrocautery (burning)
\item Ligation (tying off the base)
\item No treatment needed for asymptomatic tags
\end{itemize}

\section*{References}
\begin{thebibliography}{99}
\bibitem{jaad-skin-tags} Higgins JC, et al. "Skin tags: clinical features and management." \emph{J Am Acad Dermatol}. 2024;90(4):722-732.
\bibitem{uptodate-skin-tags} Goldstein BG, et al. Benign cutaneous lesions: skin tags. In: UpToDate, Post TW (Ed), UpToDate, Waltham, MA; 2025.
\bibitem{dermatol-surg-tags} Wollina U, et al. "Treatment of skin tags: a systematic review." \emph{Dermatol Surg}. 2023;49(6):548-556.
\bibitem{bmj-tags} Oakley A. "Skin tags." \emph{BMJ}. 2024;384:q156.
\bibitem{harrison-skin} Jameson JL, et al. \emph{Harrison\'s Principles of Internal Medicine}. 21st ed. McGraw-Hill; 2022. Chapter 62: Skin Disorders.
\end{thebibliography}
\clearpage
